\documentclass[a4paper,10pt]{article}

\usepackage[utf8]{inputenc}
\usepackage[T1]{fontenc}
\usepackage[french]{babel}
\usepackage[margin=2cm, top=1.8cm, bottom=1.8cm]{geometry}
\usepackage{titlesec}
\usepackage{enumitem}
\usepackage{hyperref}
\usepackage{xcolor}
\usepackage{fontawesome5}
\usepackage{array}
\usepackage{tabularx}
\usepackage{parskip}

% Colors
\definecolor{accent}{RGB}{80, 100, 200}
\definecolor{darkgray}{RGB}{50, 50, 50}
\definecolor{lightgray}{RGB}{130, 130, 130}

% Hyperref setup
\hypersetup{
    colorlinks=true,
    urlcolor=accent,
    linkcolor=accent
}

% Section formatting
\titleformat{\section}
  {\large\bfseries\color{darkgray}}
  {}{0em}{}
  [\vspace{0.3em}\color{lightgray}\hrule\vspace{0.4em}]

\titlespacing{\section}{0pt}{1.2em}{0.5em}

\pagestyle{empty}

\setlist[itemize]{noitemsep, topsep=2pt, leftmargin=1.2em}

% Custom commands
\newcommand{\entry}[4]{%
  \noindent\textbf{#1} \hfill \textcolor{accent}{#2}\\
  \textit{\textcolor{lightgray}{#3}}%
  \ifx&#4&\else\\\vspace{0.2em}#4\fi
  \vspace{0.5em}
}

\newcommand{\entrylist}[5]{%
  \noindent\textbf{#1} \hfill \textcolor{accent}{#2}\\
  \textit{\textcolor{lightgray}{#3}}%
  \ifx&#4&\else\\\vspace{0.2em}#4\fi
  \vspace{0.3em}
  \begin{itemize}#5\end{itemize}
  \vspace{0.3em}
}

\newcommand{\langentry}[3]{%
  \noindent\textbf{#1} — #2%
  \ifx&#3&\else\\\textcolor{lightgray}{\small #3}\fi
  \vspace{0.4em}
}

\newcommand{\skillgroup}[2]{%
  \noindent\textbf{#1 :} #2\par\vspace{0.35em}
}

%--------------------------
\begin{document}

%--- HEADER ---%
\begin{center}
  {\Huge\bfseries\color{darkgray} Mattia Giovannini}\\[0.4em]
  {\large\color{accent} Étudiant en ingénierie informatique -- IA et analyse de données}\\[0.7em]
  \textcolor{lightgray}{
    \faEnvelope\ \href{mailto:mattia.giovannini6@studio.unibo.it}{mattia.giovannini6@studio.unibo.it}
    \quad\textbullet\quad
    \faGithub\ \href{https://github.com/frigori}{github.com/frigori}
  }
\end{center}

\vspace{0.5em}

%--- PROFIL ---%
\section{Profil}

Étudiant en dernière année de licence en ingénierie informatique à l'Université de Bologne, avec un intérêt marqué pour l'intelligence artificielle, le machine learning et l'analyse de données. Ma formation m'a donné des bases solides en algorithmique, statistiques, systèmes et génie logiciel, que j'ai appliquées dans un contexte de recherche lors d'un semestre Erasmus en France. Je souhaite poursuivre un master en informatique / intelligence artificielle en France, idéalement en alternance, afin de développer des solutions intelligentes appliquées à des problèmes concrets.

%--- FORMATION ---%
\section{Formation}

\entrylist{Licence en ingénierie informatique}{2021 -- juillet 2026 (prévu)}{Université de Bologne}{Moyenne pondérée : 25,67/30}{
  \item Cours pertinents : analyse mathématique, algèbre linéaire et géométrie, probabilités et statistiques, algorithmique, systèmes d'exploitation, réseaux informatiques, systèmes d'information, génie logiciel, technologies web, architecture des ordinateurs.
  \item \textbf{Mémoire :} \textit{Détection du trouble du spectre de l'autisme à partir de données d'eye-tracking par machine learning} — classification et analyse de signaux appliquées à des enregistrements cliniques d'oculométrie.
}

\entry{Diplôme de lycée scientifique}{2016 -- 2021}{Liceo Augusto Righi, Bologne}{}

%--- EXPERIENCE INTERNATIONALE ---%
\section{Projet d'études Erasmus}

\entrylist{Expérience Erasmus}{2025}{Université de Tours, France}{}{
  \item Réalisation d'un projet sur le \textbf{machine learning appliqué à l'analyse du trouble du spectre de l'autisme (TSA)}.
  \item Application de techniques de classification et d'extraction de caractéristiques à des données d'eye-tracking.
  \item Expérience dans un environnement académique international et francophone.
  \item Renforcement des compétences en analyse de données, Python et méthodologie scientifique.
}

%--- PROJETS ---%
\section{Expérience projet}

\entrylist{BOOM GEN AI Challenge 2026 — VEM Sistemi}{2026}{Équipe projet GenAI}{}{
  \item Contribution au développement d'un \textbf{assistant GenAI} pour l'analyse automatisée d'appels d'offres et l'aide à la décision \textbf{GO/NO-GO}.
  \item Conception et implémentation d'un \textbf{module de scoring en Python} pour évaluer l'alignement entre les exigences d'un appel d'offres et les capacités de l'entreprise.
  \item Expérimentation de techniques de \textbf{prompt engineering} et de services \textbf{Azure OpenAI} pour l'extraction, la classification et l'enrichissement des exigences.
  \item Analyse d'un pipeline de \textbf{parsing documentaire} pour normaliser des documents non structurés.
}

%--- COMPETENCES ---%
\section{Compétences techniques}

\skillgroup{Langages}{C, Java, Python, JavaScript, HTML/CSS, SQL, Bash}
\skillgroup{Machine learning \& données}{scikit-learn, NumPy, pandas, Matplotlib, Jupyter Notebook}
\skillgroup{Systèmes \& réseaux}{Linux, TCP/IP, programmation socket, bases de données relationnelles}
\skillgroup{Outils de développement}{Git, \LaTeX, MATLAB, utilisation de Docker \& Podman}
\skillgroup{Centres d'intérêt}{Machine learning, analyse de données, IA générative, conteneurisation, automatisation système}

%--- AUTRES EXPERIENCES ---%
\section{Autres expériences}

\entrylist{Webmaster et rédacteur étudiant}{2020 -- 2021}{Liceo Augusto Righi, Bologne}{}{
  \item Conception et gestion de plateformes étudiantes (\textit{Righi Network} et \textit{Il Sottobanco}), au service de la communication numérique et du journal de l'école.
  \item Développement de compétences en technologies web, gestion de contenu et travail en équipe.
}

\entry{Animateur bénévole en centre d'été}{2016 -- 2020}{Estate Ragazzi, Bologne}{Gestion de groupe, travail en équipe, communication et responsabilité dans des activités éducatives.}

%--- LANGUES ---%
\section{Langues}

\langentry{Italien}{Langue maternelle}{}

\langentry{Anglais}{C1}{Duolingo English Test : 130/160 (C1) · Certification CLA C1, Université de Bologne · Cambridge First Certificate B2 (178/190)}

\langentry{Français}{B1/B2}{Certification CLA B1, Université de Bologne · cours B2/C1 suivis à l'Université de Tours}

\end{document}
